\documentclass[a4paper,12pt]{ctexart}
\usepackage{../../teaching-style}

\begin{document}
\pagestyle{empty}

\maintitle{概率论与数理统计}
\begin{center}
  \large 高学段 \quad ★ 标记为高考不考内容
\end{center}
\vspace{0.3cm}

\chaptertitle{第一章\quad 排列组合与二项式定理}

\sectiontitle{1.1\quad 排列组合回顾}

\knowbox{
$A_n^m=n(n-1)\cdots(n-m+1)=\dfrac{n!}{(n-m)!}$。\\
$C_n^m=\dfrac{A_n^m}{A_m^m}=\dfrac{n!}{m!(n-m)!}$。\\
性质：$C_n^m=C_n^{n-m}$，$C_{n+1}^m=C_n^m+C_n^{m-1}$（杨辉恒等式）。}

\sectiontitle{1.2\quad 二项式定理}

\knowbox{
$(a+b)^n=\sum_{k=0}^n C_n^k a^{n-k}b^k$。\\
通项：$T_{k+1}=C_n^k a^{n-k}b^k$。\\
性质：二项式系数和$C_n^0+C_n^1+\cdots+C_n^n=2^n$。\\
奇数项系数和$=$偶数项系数和$=2^{n-1}$。}

\sectiontitle{习题1}

\begin{exercise}
\item 求$(x+2)^5$展开式中$x^3$的系数。
\item 求$(2x-\dfrac1x)^6$的常数项。
\item $(1+x)^n$展开式中$x^3$的系数是$x^2$系数的$2$倍，求$n$。
\end{exercise}

\newpage

\chaptertitle{第二章\quad 概率进阶}

\sectiontitle{2.1\quad 条件概率与乘法公式}

\knowbox{
$P(B|A)=\dfrac{P(AB)}{P(A)}$（$P(A)>0$）。\\
乘法公式：$P(AB)=P(A)P(B|A)=P(B)P(A|B)$。\\
事件的独立性：$P(AB)=P(A)P(B)\Leftrightarrow A,B$独立。}

\sectiontitle{2.2\quad 全概率公式与贝叶斯公式 ★}

\knowbox{
\textbf{全概率公式}：若$B_1,B_2,\dots,B_n$是完备事件组，则$P(A)=\sum_{i=1}^n P(B_i)P(A|B_i)$。\\
\textbf{贝叶斯公式}：$P(B_i|A)=\dfrac{P(B_i)P(A|B_i)}{\sum_{j=1}^n P(B_j)P(A|B_j)}$。\\
意义：用结果反推原因的概率（后验概率）。}

\sectiontitle{习题2}

\begin{exercise}
\item 袋中有$3$红$2$白球，先后取两球（不放回），求第二次取到红球的概率。
\item 某产品合格率$90\%$，合格品中一等品率$80\%$，求产品是一等品的概率。
\item ★ 两台机器生产同种产品，甲产量$60\%$，合格率$95\%$；乙产量$40\%$，合格率$90\%$。\\
  求任取一件是合格品的概率。若取到一件合格品，求它是甲厂生产的概率。
\end{exercise}

\newpage

\chaptertitle{第三章\quad 随机变量与分布}

\sectiontitle{3.1\quad 随机变量与分布列}

\knowbox{
\textbf{离散型随机变量}$X$：取值有限或可列个。\\
分布列：$P(X=x_i)=p_i$，满足$p_i\ge0$，$\sum p_i=1$。\\
\textbf{分布函数}：$F(x)=P(X\le x)$，单调不减，右连续。}

\sectiontitle{3.2\quad 常见离散型分布}

\knowbox{
\textbf{二项分布}$X\sim B(n,p)$：$P(X=k)=C_n^k p^k(1-p)^{n-k}$（$n$次独立重复试验）。\\
期望$E(X)=np$，方差$D(X)=np(1-p)$。\\
\textbf{超几何分布}$X\sim H(N,M,n)$：$P(X=k)=\dfrac{C_M^k C_{N-M}^{n-k}}{C_N^n}$。\\
\textbf{泊松分布}★ $X\sim P(\lambda)$：$P(X=k)=\dfrac{\lambda^k}{k!}e^{-\lambda}$（$k=0,1,2,\dots$）。}

\sectiontitle{3.3\quad 正态分布 ★}

\knowbox{
连续型随机变量，密度函数$f(x)=\dfrac1{\sqrt{2\pi}\sigma}e^{-\frac{(x-\mu)^2}{2\sigma^2}}$。\\
$X\sim N(\mu,\sigma^2)$，$E(X)=\mu$，$D(X)=\sigma^2$。\\
标准正态$N(0,1)$：$\Phi(x)=\dfrac1{\sqrt{2\pi}}\int_{-\infty}^x e^{-\frac{t^2}{2}}dt$。\\
$3\sigma$原则：$P(|X-\mu|<\sigma)\approx68.3\%$，$2\sigma\approx95.4\%$，$3\sigma\approx99.7\%$。}

\sectiontitle{习题3}

\begin{exercise}
\item 掷硬币$5$次，求正面朝上次数的分布列和期望。
\item 某射击手命中率$0.8$，射击$10$次，求至少命中$8$次的概率。
\item ★ 某地区成年男性身高$X\sim N(170,10^2)$，求身高在$160$到$180$之间的概率。
\end{exercise}

\newpage

\chaptertitle{第四章\quad 期望与方差}

\sectiontitle{4.1\quad 数学期望}

\knowbox{
离散型：$E(X)=\sum x_i p_i$。\\
连续型★：$E(X)=\int_{-\infty}^{\infty} x f(x)\,dx$。\\
性质：$E(aX+b)=aE(X)+b$；$E(X+Y)=E(X)+E(Y)$。\\
若$X,Y$独立，$E(XY)=E(X)E(Y)$。}

\sectiontitle{4.2\quad 方差}

\knowbox{
$D(X)=E[(X-E(X))^2]=E(X^2)-[E(X)]^2$。\\
性质：$D(aX+b)=a^2 D(X)$；若$X,Y$独立，$D(X\pm Y)=D(X)+D(Y)$。\\
标准差$\sigma(X)=\sqrt{D(X)}$。}

\sectiontitle{4.3\quad 大数定律与中心极限定理 ★}

\knowbox{
\textbf{切比雪夫不等式}：$P(|X-\mu|\ge\varepsilon)\le\dfrac{\sigma^2}{\varepsilon^2}$。\\
\textbf{大数定律}：样本均值依概率收敛于总体均值。\\
\textbf{中心极限定理}：独立同分布随机变量的和近似服从正态分布。\\
重要性：解释了正态分布在自然界中普遍存在的原因。}

\sectiontitle{习题4}

\begin{exercise}
\item 已知$X$分布列：$P(X=0)=0.2$，$P(X=1)=0.5$，$P(X=2)=0.3$，求$E(X)$和$D(X)$。
\item $X\sim B(10,0.2)$，求$E(X)$和$D(X)$。
\item 已知$E(X)=3$，$D(X)=4$，求$E(2X-1)$和$D(2X-1)$。
\end{exercise}

\newpage

\chaptertitle{第五章\quad 数理统计初步}

\sectiontitle{5.1\quad 抽样方法与统计量}

\knowbox{
\textbf{简单随机抽样}：每个个体被抽到的概率相等。\\
\textbf{分层抽样}：按比例从各层抽取。\\
\textbf{系统抽样}：等距抽取。\\
\textbf{样本均值}$\bar X=\dfrac1n\sum_{i=1}^n X_i$。\\
\textbf{样本方差}$S^2=\dfrac1{n-1}\sum_{i=1}^n (X_i-\bar X)^2$（无偏估计）。}

\sectiontitle{5.2\quad 参数估计 ★}

\knowbox{
\textbf{点估计}：用样本统计量估计总体参数。\\
$\hat\mu=\bar X$，$\hat{\sigma^2}=S^2$。\\
\textbf{区间估计}：$\bar X\pm z_{\alpha/2}\dfrac{\sigma}{\sqrt{n}}$（$\sigma$已知时总体均值的置信区间）。\\
置信水平越高，区间越宽。}

\sectiontitle{5.3\quad 假设检验与回归分析 ★}

\knowbox{
\textbf{假设检验}：提出假设$H_0$，计算检验统计量，根据$p$值判断是否拒绝$H_0$。\\
两类错误：弃真（$\alpha$）和纳伪（$\beta$）。\\
\textbf{线性回归}：$y=\beta_0+\beta_1 x+\varepsilon$。\\
最小二乘法：使$\sum(y_i-\hat y_i)^2$最小。\\
相关系数$r$：衡量$x$与$y$线性相关程度（$|r|\le1$）。\\
$r$接近$\pm1$表示强线性相关，接近$0$表示弱相关。}

\sectiontitle{习题5}

\begin{exercise}
\item 某样本数据：$12,15,14,13,16$，求样本均值和样本方差。
\item 从$500$名学生中抽取$50$名调查，采用什么抽样方法比较合理？
\item ★ 某厂产品寿命$\sim N(\mu,\sigma^2)$，抽样$n=100$，$\bar x=1050$小时，$s=40$小时，求$\mu$的$95\%$置信区间（$z_{0.025}\approx1.96$）。
\item ★ 给出$5$组数据$(x,y)$：$(1,2),(2,4),(3,5),(4,7),(5,8)$，求回归直线方程。
\end{exercise}

\end{document}
